\chapter{Looking Back on the Year}

% ================= PERCAKAPAN 1-2 =================
\begin{convobox}{1}
    \noindent
    \jpkana{こんにちは。}{Konnichiwa.}

    \vspace{1.5em}
    \noindent{\Large \color{articolor} \textbf{Arti:} Halo.}
\end{convobox}

\begin{convobox}{2}
    \noindent
    \jpword{ことし}{今年}{Kotoshi} \jpkana{も}{mo} 
    \jpkana{あと}{ato} \jpword{すこ}{少}{suko}\jpkana{しで}{shi de} 
    \jpword{お}{終}{o}\jpkana{わりですね。}{wari desu ne.}

    \vspace{1.5em}
    \noindent{\Large \color{articolor} \textbf{Arti:} Tahun ini juga tinggal sedikit lagi selesai ya.}
\end{convobox}

\begin{lessonbox}{Catatan Belajar: Tinggal Sedikit Lagi!}
    Halo adik-adik! Kalau kamu sedang menunggu sesuatu yang hampir selesai, kamu bisa pakai kata ajaib \textbf{あと少しで (Ato sukoshi de)} yang artinya "Tinggal sedikit lagi". Contohnya: \textit{Ato sukoshi de liburan!} Wah, asik banget!
\end{lessonbox}

% ================= PERCAKAPAN 3-4 =================
\begin{convobox}{3}
    \noindent
    \jpkana{ゆっくり}{Yukkuri} 
    \jpword{ちゃ}{お茶}{o-cha} \jpkana{でも}{demo} 
    \jpword{の}{飲}{no}\jpkana{みながら、}{mi nagara,}

    \vspace{1.5em}
    \noindent{\Large \color{articolor} \textbf{Arti:} Sambil bersantai meminum teh atau semacamnya,}
\end{convobox}

\begin{convobox}{4}
    \noindent
    \jpword{いっしょ}{一緒}{Issho} \jpkana{に}{ni} 
    \jpkana{この}{kono} 
    \jpword{いちねん}{1年}{ichi-nen} \jpkana{を}{o} 
    \jpword{ふ}{振}{fu}\jpkana{り}{ri}\jpword{かえ}{返}{kae}\jpkana{りましょう。}{rimashou.}

    \vspace{1.5em}
    \noindent{\Large \color{articolor} \textbf{Arti:} mari kita kenang (ingat kembali) 1 tahun ini bersama-sama.}
\end{convobox}

\begin{lessonbox}{Catatan Belajar: Mari Kita...!}
    Kalau kamu mau mengajak temanmu melakukan sesuatu, ubah akhiran kata kerjanya menjadi \textbf{〜ましょう (〜mashou)}. Artinya adalah "Mari kita\dots". Contoh: \textit{Furikaerimashou!} (Mari kita ingat-ingat!).
\end{lessonbox}

% ================= PERCAKAPAN 5-6 =================
\begin{convobox}{5}
    \noindent
    \jpword{ことし}{今年}{Kotoshi} \jpkana{は、}{wa,} 
    \jpword{たの}{楽}{tano}\jpkana{しい}{shii} 
    \jpword{いちねん}{1年}{ichi-nen} \jpkana{でしたか？}{deshita ka?}

    \vspace{1.5em}
    \noindent{\Large \color{articolor} \textbf{Arti:} Apakah tahun ini adalah tahun yang menyenangkan?}
\end{convobox}

\begin{convobox}{6}
    \noindent
    \jpkana{それとも、}{Sore tomo,} 
    \jpword{たいへん}{大変}{taihen} \jpkana{な}{na} 
    \jpword{いちねん}{1年}{ichi-nen} \jpkana{でしたか？}{deshita ka?}

    \vspace{1.5em}
    \noindent{\Large \color{articolor} \textbf{Arti:} Atau, apakah tahun yang berat?}
\end{convobox}

\begin{lessonbox}{Catatan Belajar: Atau yang Mana?}
    Saat kamu memberikan 2 pilihan pertanyaan, gunakan kata penyambung \textbf{それとも (Soretomo)} yang artinya "Atau". Menyenangkan? \textit{Soretomo} (Atau) berat?
\end{lessonbox}

% ================= PERCAKAPAN 7-8 =================
\begin{convobox}{7}
    \noindent
    \jpkana{どちらかというと、}{Dochira ka to iu to,} 
    \jpword{たいへん}{大変}{taihen} \jpkana{な}{na} 
    \jpword{いちねん}{1年}{ichi-nen} \jpkana{でした。}{deshita.}

    \vspace{1.5em}
    \noindent{\Large \color{articolor} \textbf{Arti:} Kalau boleh dibilang, ini tahun yang berat.}
\end{convobox}

\begin{convobox}{8}
    \noindent
    \jpkana{そうですか。}{Sou desu ka.}

    \vspace{1.5em}
    \noindent{\Large \color{articolor} \textbf{Arti:} Oh begitu.}
\end{convobox}

\begin{lessonbox}{Catatan Belajar: "Kalau Disuruh Milih..."}
    Nah, \textbf{どちらかというと (Dochira ka to iu to)} ini adalah kalimat gaul orang Jepang kalau mereka ragu-ragu tapi harus memilih! Artinya kurang lebih: "Kalau disuruh milih sih\dots" atau "Bisa dibilang sih\dots".
\end{lessonbox}

% ================= PERCAKAPAN 9-10 =================
\begin{convobox}{9}
    \noindent
    \jpword{ことし}{今年}{Kotoshi}\jpkana{、}{,} 
    \jpword{いちばん}{一番}{ichiban} 
    \jpword{たの}{楽}{tano}\jpkana{しかった}{shikatta} 
    \jpkana{こと}{koto} \jpkana{は}{wa} 
    \jpword{なん}{何}{nan} \jpkana{でしたか？}{deshita ka?}

    \vspace{1.5em}
    \noindent{\Large \color{articolor} \textbf{Arti:} Tahun ini, hal apa yang paling menyenangkan?}
\end{convobox}

\begin{convobox}{10}
    \noindent
    \jpword{なつ}{夏}{Natsu} \jpkana{に、}{ni,} 
    \jpword{かぞく}{家族}{kazoku} \jpkana{で}{de} 
    \jpkana{キャンプ}{kyanpu} \jpkana{に}{ni} 
    \jpword{い}{行}{i}\jpkana{った}{tta} 
    \jpkana{こと}{koto} \jpkana{です。}{desu.}

    \vspace{1.5em}
    \noindent{\Large \color{articolor} \textbf{Arti:} (Hal yang menyenangkan adalah) Pergi berkemah bersama keluarga di musim panas.}
\end{convobox}

\begin{lessonbox}{Catatan Belajar: Paling Juara!}
    Adik-adik tahu juara nomor 1? Di Jepang, kata untuk "Paling" atau "Nomor Satu" adalah \textbf{一番 (Ichiban)}! Jadi \textit{Ichiban tanoshikatta} artinya "Yaaang paling menyenangkan!".
\end{lessonbox}

% ================= PERCAKAPAN 11-12 =================
\begin{convobox}{11}
    \noindent
    \jpkana{では、}{Dewa,} 
    \jpword{ことし}{今年}{kotoshi}\jpkana{、}{,} 
    \jpword{いちばん}{一番}{ichiban} 
    \jpword{たいへん}{大変}{taihen} \jpkana{だった}{datta} 
    \jpkana{こと}{koto} \jpkana{は}{wa} 
    \jpword{なん}{何}{nan} \jpkana{でしたか？}{deshita ka?}

    \vspace{1.5em}
    \noindent{\Large \color{articolor} \textbf{Arti:} Kalau begitu, tahun ini, hal apa yang paling berat?}
\end{convobox}

\begin{convobox}{12}
    \noindent
    \jpkana{バイク}{Baiku} \jpkana{で}{de} 
    \jpword{ころ}{転}{koro}\jpkana{んで、}{nde,} 
    \jpword{こっせつ}{骨折}{kossetsu} \jpkana{した}{shita} 
    \jpkana{こと}{koto} \jpkana{です。}{desu.}

    \vspace{1.5em}
    \noindent{\Large \color{articolor} \textbf{Arti:} Jatuh dari (menggunakan) motor dan patah tulang.}
\end{convobox}

\begin{lessonbox}{Catatan Belajar: Naik Apa?}
    Coba lihat partikel \textbf{で (De)} setelah kata motor (\textit{Baiku}). Di bahasa Jepang, kalau mau bilang kita pergi "Naik" atau "Menggunakan" suatu alat, kita pakai partikel ini! \textit{Baiku de} = Dengan motor. Jangan sampai jatuh ya!
\end{lessonbox}

% ================= PERCAKAPAN 13-14 =================
\begin{convobox}{13}
    \noindent
    \jpkana{それは}{Sore wa} 
    \jpword{たいへん}{大変}{taihen} \jpkana{でしたね。}{deshita ne.}

    \vspace{1.5em}
    \noindent{\Large \color{articolor} \textbf{Arti:} Itu pasti sangat berat (menyusahkan) ya.}
\end{convobox}

\begin{convobox}{14}
    \noindent
    \jpword{ことし}{今年}{Kotoshi}\jpkana{、}{,} 
    \jpword{いちばん}{一番}{ichiban} 
    \jpword{か}{買}{ka}\jpkana{って}{tte} 
    \jpkana{よかった}{yokatta} 
    \jpword{もの}{物}{mono} \jpkana{は}{wa} 
    \jpword{なん}{何}{nan} \jpkana{でしたか？}{deshita ka?}

    \vspace{1.5em}
    \noindent{\Large \color{articolor} \textbf{Arti:} Tahun ini, barang apa yang paling kamu syukuri karena telah membelinya?}
\end{convobox}

\begin{lessonbox}{Catatan Belajar: Syukurlah Sudah...!}
    Kalau kamu senang melakukan sesuatu, kamu bisa bilang "Syukurlah sudah...". Caranya gampang! Tambahkan \textbf{〜てよかった (〜te yokatta)} di belakang kata kerja. Contohnya \textit{Katte yokatta} = Syukurlah sudah beli (Tidak menyesal beli).
\end{lessonbox}

% ================= PERCAKAPAN 15-16 =================
\begin{convobox}{15}
    \noindent
    \jpkana{ノイズキャンセリング}{Noizu kyanseringu} \jpkana{が}{ga} 
    \jpkana{ついた}{tsuita} 
    \jpkana{イヤホン}{iyahon} \jpkana{です。}{desu.}

    \vspace{1.5em}
    \noindent{\Large \color{articolor} \textbf{Arti:} Earphone yang dilengkapi noise-canceling.}
\end{convobox}

\begin{convobox}{16}
    \noindent
    \jpword{ことし}{今年}{Kotoshi}\jpkana{、}{,} 
    \jpword{いちばん}{一番}{ichiban} 
    \jpkana{がんばった}{ganbatta} 
    \jpkana{こと}{koto} \jpkana{は}{wa} 
    \jpword{なん}{何}{nan} \jpkana{でしたか？}{deshita ka?}

    \vspace{1.5em}
    \noindent{\Large \color{articolor} \textbf{Arti:} Tahun ini, hal apa yang paling kamu usahakan (perjuangkan)?}
\end{convobox}

\begin{lessonbox}{Catatan Belajar: Berjuang Keras!}
    Kamu sering dengar kata "Ganbatte!" kan? Nah, kalau perjuangannya SUDAH selesai atau berlalu, kita ubah menjadi bentuk lampau yaitu \textbf{がんばった (Ganbatta)}! Kamu juga \textit{ganbatta} lho belajar bahasa Jepang tahun ini!
\end{lessonbox}

% ================= PERCAKAPAN 17-18 =================
\begin{convobox}{17}
    \noindent
    \jpkana{アルバイト}{Arubaito} \jpkana{です。}{desu.}

    \vspace{1.5em}
    \noindent{\Large \color{articolor} \textbf{Arti:} Kerja paruh waktu.}
\end{convobox}

\begin{convobox}{18}
    \noindent
    \jpword{ことし}{今年}{Kotoshi}\jpkana{、}{,} 
    \jpword{いちばん}{一番}{ichiban} \jpkana{の}{no} 
    \jpword{まな}{学}{mana}\jpkana{びは}{bi wa} 
    \jpword{なん}{何}{nan} \jpkana{でしたか？}{deshita ka?}

    \vspace{1.5em}
    \noindent{\Large \color{articolor} \textbf{Arti:} Tahun ini, apa pelajaran (hikmah) paling berharga?}
\end{convobox}

\begin{lessonbox}{Catatan Belajar: Pelajaran Berharga}
    Dari kata \textit{Manabu} (Belajar), kita bisa buat kata benda \textbf{学び (Manabi)} yang berarti "Pelajaran hidup" atau "Hikmah". Apa nih \textit{Manabi} kamu tahun ini?
\end{lessonbox}

% ================= PERCAKAPAN 19-20 =================
\begin{convobox}{19}
    \noindent
    \jpword{はたら}{働}{Hatara}\jpkana{きすぎて、}{ki-sugite,} 
    \jpword{たいちょう}{体調}{taichou} \jpkana{を}{o} 
    \jpword{くず}{崩}{kuzu}\jpkana{してしまった}{shite shimatta} 
    \jpword{とき}{時}{toki}\jpkana{、}{,}

    \vspace{1.5em}
    \noindent{\Large \color{articolor} \textbf{Arti:} Saat tubuhku jatuh sakit karena terlalu banyak bekerja,}
\end{convobox}

\begin{convobox}{20}
    \noindent
    \jpword{けんこう}{健康}{Kenkou} \jpkana{が}{ga} 
    \jpword{いちばん}{一番}{ichiban} 
    \jpword{だいじ}{大事}{daiji} \jpkana{だと}{da to} 
    \jpword{まな}{学}{mana}\jpkana{びました。}{bimashita.}

    \vspace{1.5em}
    \noindent{\Large \color{articolor} \textbf{Arti:} aku belajar bahwa kesehatan adalah yang paling penting.}
\end{convobox}

\begin{lessonbox}{Catatan Belajar: Jangan Terlalu Banyak!}
    Kalau kamu melakukan sesuatu secara berlebihan, tambahkan \textbf{〜すぎる (〜sugiru)} di akhir kata kerja! Makan terlalu banyak = \textit{Tabe-sugiru}. Kalau di sini, \textit{Hataraki-sugite} = Terlalu banyak bekerja. Ingat, \textit{kenkou} (kesehatan) itu penting lho!
\end{lessonbox}

% ================= PERCAKAPAN 21-22 =================
\begin{convobox}{21}
    \noindent
    \jpword{らいねん}{来年}{Rainen}\jpkana{、}{,} 
    \jpword{ことし}{今年}{kotoshi} \jpkana{よりも}{yori mo} 
    \jpkana{がんばりたい}{ganbaritai} 
    \jpkana{こと}{koto} \jpkana{は}{wa} 
    \jpword{なん}{何}{nan} \jpkana{ですか？}{desu ka?}

    \vspace{1.5em}
    \noindent{\Large \color{articolor} \textbf{Arti:} Tahun depan, hal apa yang ingin kamu usahakan lebih daripada tahun ini?}
\end{convobox}

\begin{convobox}{22}
    \noindent
    \jpword{らいねん}{来年}{Rainen} \jpkana{は、}{wa,} 
    \jpword{ことし}{今年}{kotoshi} \jpkana{よりも}{yori mo} 
    \jpkana{たくさん}{takusan} 
    \jpkana{ジム}{jimu} \jpkana{に}{ni} 
    \jpword{かよ}{通}{kayo}\jpkana{いたいです。}{itai desu.}

    \vspace{1.5em}
    \noindent{\Large \color{articolor} \textbf{Arti:} Tahun depan, aku ingin lebih banyak (sering) pergi ke gym dibandingkan tahun ini.}
\end{convobox}

\begin{lessonbox}{Catatan Belajar: "Daripada..."}
    Kamu mau membandingkan 2 hal? Pakai saja grammar \textbf{〜よりも (〜yori mo)}. Artinya "Dibandingkan / Daripada". Contoh: \textit{Kotoshi yori mo} = Daripada tahun ini.
\end{lessonbox}

% ================= PERCAKAPAN 23-24 =================
\begin{convobox}{23}
    \noindent
    \jpkana{では、}{Dewa,} 
    \jpword{らいねん}{来年}{rainen}\jpkana{、}{,} 
    \jpword{あたら}{新}{atara}\jpkana{しく}{shiku} 
    \jpword{ちょうせん}{挑戦}{chousen} \jpkana{してみたい}{shite mitai} 
    \jpkana{こと}{koto} \jpkana{は}{wa}

    \vspace{1.5em}
    \noindent{\Large \color{articolor} \textbf{Arti:} Kalau begitu, tahun depan, hal apa yang ingin kamu coba tantang (lakukan) yang baru}
\end{convobox}

\begin{convobox}{24}
    \noindent
    \jpword{なに}{何}{nani}\jpkana{か}{ka} 
    \jpkana{ありますか？}{arimasu ka?}

    \vspace{1.5em}
    \noindent{\Large \color{articolor} \textbf{Arti:} apakah ada?}
\end{convobox}

\begin{lessonbox}{Catatan Belajar: "Aku Ingin Mencoba!"}
    Kalau kamu penasaran dan ingin mencoba sesuatu yang baru, gunakan rumus \textbf{〜てみたい (〜te mitai)}. Berasal dari kata "Melihat", di sini artinya "Ingin mencoba melakukan". Keren kan?
\end{lessonbox}

% ================= PERCAKAPAN 25-26 =================
\begin{convobox}{25}
    \noindent
    \jpword{からて}{空手}{Karate} \jpkana{を}{o} 
    \jpword{なら}{習}{nara}\jpkana{ってみたいです。}{tte mitai desu.}

    \vspace{1.5em}
    \noindent{\Large \color{articolor} \textbf{Arti:} Aku ingin mencoba belajar Karate.}
\end{convobox}

\begin{convobox}{26}
    \noindent
    \jpword{ことし}{今年}{Kotoshi} 
    \jpword{いちねん}{1年}{ichi-nen}\jpkana{、}{,} 
    \jpword{つら}{辛}{tsura}\jpkana{かった}{katta} 
    \jpkana{こと}{koto} \jpkana{や}{ya}

    \vspace{1.5em}
    \noindent{\Large \color{articolor} \textbf{Arti:} Selama 1 tahun ini, hal-hal yang menyakitkan (sulit) dan}
\end{convobox}

\begin{lessonbox}{Catatan Belajar: Menyebutkan Banyak Hal!}
    Partikel \textbf{や (Ya)} dipakai kalau kamu mau bilang "Dan" atau "Atau" saat menyebutkan sebagian contoh! Misalnya, beli Apel \textit{ya} Jeruk \textit{ya} Pisang.
\end{lessonbox}

% ================= PERCAKAPAN 27-28 =================
\begin{convobox}{27}
    \noindent
    \jpword{たいへん}{大変}{taihen} \jpkana{だった}{datta} 
    \jpkana{こと}{koto} \jpkana{も}{mo} 
    \jpkana{あったと}{atta to} 
    \jpword{おも}{思}{omo}\jpkana{います。}{imasu.}

    \vspace{1.5em}
    \noindent{\Large \color{articolor} \textbf{Arti:} hal-hal yang berat, aku pikir pasti (pernah) ada.}
\end{convobox}

\begin{convobox}{28}
    \noindent
    \jpkana{でも、}{Demo,} 
    \jpkana{もうすぐ}{mou sugu} 
    \jpword{あたら}{新}{atara}\jpkana{しい}{shii} 
    \jpword{とし}{年}{toshi} \jpkana{が}{ga} 
    \jpword{はじ}{始}{haji}\jpkana{まります。}{marimasu.}

    \vspace{1.5em}
    \noindent{\Large \color{articolor} \textbf{Arti:} Tapi, sebentar lagi tahun yang baru akan dimulai.}
\end{convobox}

\begin{lessonbox}{Catatan Belajar: Sebentar Lagi!}
    Wah, kata \textbf{もうすぐ (Mou sugu)} sangat berguna loh! Artinya "Sebentar lagi". \textit{Mou sugu} makan malam! \textit{Mou sugu} tahun baru!
\end{lessonbox}

% ================= PERCAKAPAN 29-31 =================
\begin{convobox}{29}
    \noindent
    \jpword{らいねん}{来年}{Rainen} \jpkana{は、}{wa,} 
    \jpword{ことし}{今年}{kotoshi} \jpkana{よりも}{yori mo} 
    \jpkana{もっと}{motto} 
    \jpword{すてき}{素敵}{suteki} \jpkana{な}{na} 
    \jpkana{こと}{koto} \jpkana{が}{ga}

    \vspace{1.5em}
    \noindent{\Large \color{articolor} \textbf{Arti:} Di tahun depan, semoga ada lebih banyak hal yang indah}
\end{convobox}

\begin{convobox}{30}
    \noindent
    \jpkana{たくさん}{takusan} 
    \jpkana{ありますように。}{arimasu you ni.}

    \vspace{1.5em}
    \noindent{\Large \color{articolor} \textbf{Arti:} (lanjutan) daripada tahun ini.}
\end{convobox}

\begin{convobox}{31}
    \noindent
    \jpkana{それでは、}{Sore dewa,} 
    \jpkana{よいお}{yoi o}\jpword{とし}{年}{toshi} \jpkana{を。}{o.}

    \vspace{1.5em}
    \noindent{\Large \color{articolor} \textbf{Arti:} Kalau begitu, selamat menyambut tahun baru.}
\end{convobox}

\begin{lessonbox}{Catatan Belajar: Doa dan Harapan Tahun Baru}
    Bila kamu mau mendoakan hal baik, akhiri kalimatmu dengan \textbf{〜ように (〜you ni)}! Artinya "Semoga\dots". Dan yang terakhir, jangan lupa ucapkan \textbf{よいお年を (Yoi o-toshi o)} ke orang-orang sebelum pergantian tahun, ya!
\end{lessonbox}

% --- SUMMARY BOX ---
\vspace{2em}
\begin{tcolorbox}[
    colback=blue!5!white,
    colframe=blue!80!black,
    boxrule=3pt,
    arc=5mm,
    title={\huge\bfseries \sffamily 🌟 Rangkuman Bab 8 🌟},
    fonttitle=\color{white},
    attach boxed title to top center={yshift=-4mm},
    boxed title style={colback=blue!80!black, rounded corners},
    left=5mm, right=5mm, top=7mm, bottom=5mm
]
\Large \textbf{Selamat! Kamu telah menyelesaikan pelajaran tahun ini!} Banyak sekali grammar sakti yang kita pelajari di sini:
\begin{itemize}
    \item \textbf{Mengajak:} \textit{〜mashou} (Mari kita\dots)
    \item \textbf{Bersyukur:} \textit{〜te yokatta} (Syukurlah sudah\dots)
    \item \textbf{Berlebihan:} \textit{〜sugiru} (Terlalu banyak\dots)
    \item \textbf{Perbandingan:} \textit{〜yori mo} (Daripada\dots)
    \item \textbf{Keinginan baru:} \textit{〜te mitai} (Ingin mencoba\dots)
\end{itemize}
Terima kasih sudah belajar dengan rajin bersama Andi - Japanese Handbook. Tetap pertahankan \textbf{Yaruki} (semangat) kamu ya. \textit{Yoi o-toshi o!} Semoga tahun depan lebih indah!
\end{tcolorbox}

\newpage